\section{Sistema vectorial}

\definition{Vector}{PENDIENTE}

Mientras que las magnitudes escalaes solo cuentan con unidades, los vectores cuentan también con dirección y sentido.

Los vectores cuentan con ciertas características:

\begin{itemize}
    \item Sentido: Representado por la punta de la flecha.
    \item Dirección: Recta sobre la que se plantea un vector.
    \item Módulo: Longitud entre inicio y fin del vector.
    \item Amplitud: Expresión numérica de la longitud gráfica del vector.
    \item Punto de aplicación: Lugar geométrico en que inicia el vector a nivel gráfico.
    \item Nombre: Letra que acompaña al vector.
\end{itemize}

\definition{Sistema vectorial}{Se tiene un sistema vectorial cuando cuentas con un conjunto de vectores del mismo tipo.}

Los sistemas vectoriales pueden ser de distintos tipos, esto dependiendo de las características de los diferentes vectores que componen al sistema.

\begin{itemize}
    \item Colineales: Los vectores se localizan en una sola dirección.
    \item Coplanares: Sistema en que los vectores se encuentran en el mismo plano.
    \item Concurrentes: La dirección de los vectores se cruza en algún punto (punto de aplicación).
\end{itemize}

\section{Operaciones con vectores}

Componentes rectangulares:

$r^2=x^2+y^2$
$r=\sqrt{x^2+y^2}$
$Tan(α)=\frac{y}{x}$
$α=Tan^{-1}*\frac{y}{x}$
$Tan^{-1}*\frac{-y}{-x}=α+180^°$
$Tan^{-1}*\frac{-y}{x})=α+360^°$
$Tan^{-1}*\frac{y}{-x}=α+180^°$
